%!TEX root = ../sztuthesis_main.tex
% 设置中文摘要

\addcontentsline{toc}{section}{摘要}

\keywordscn{音乐情感分析；大语言模型；音频描述生成；Valence-Arousal；DEAM 数据集}
\categorycn{TP391}
\begin{abstractcn}
随着音乐内容分析与智能推荐技术的发展，如何从音乐音频中有效提取语义信息并进一步完成情感识别，成为音乐信息检索与情感计算领域的重要研究问题。传统音乐情感分析方法多依赖声学特征与机器学习模型，虽然在数值预测上具有一定稳定性，但在高层语义表达和结果可解释性方面存在不足。针对上述问题，本文围绕“基于大语言模型的音乐音频信号描述与情感标签生成”展开研究，提出并实现了一种两阶段处理流程：首先利用预训练音频理解模型对音乐音频片段生成中文文本描述，再基于大语言模型对描述文本进行情感推理，输出 Valence、Arousal 以及离散情感标签。

本文以 DEAM 数据集为实验基础，完成了音频预处理、样本构建及训练集、验证集、测试集划分等工作。在音频描述生成阶段，系统对音乐片段进行裁剪、重采样和编码处理，并通过提示词约束模型从旋律、节奏、乐器、曲风和整体听感五个维度生成描述文本。同时，本文设计了基于规则的描述质量评估方法，对生成结果的长度、维度覆盖和幻觉风险进行自动检测。在情感推理阶段，本文将音乐情感分析建模为基于描述文本的语义推理任务，利用大语言模型输出连续情感维度值与离散情感标签，并实现了批量推理、结果解析与断点续跑等工程功能。为验证所提方法的可行性，本文进一步构建了基于传统声学特征的机器学习基线，并与大模型方法进行了初步对比分析。

实验结果表明，本文构建的系统已能够打通从音乐音频输入到描述生成、再到情感标签输出的完整处理链路，生成的描述文本在小样本测试中具有较好的可用性，能够为后续情感推理提供语义支撑。对比实验表明，当前传统基线在数值预测稳定性方面仍优于大模型链路，但本文方法在音乐语义表达与中间过程可解释性方面具有一定优势。总体而言，本文验证了“音频描述中间层 + 大语言模型情感推理”方案在音乐情感分析任务中的可行性，也指出了当前系统在实验规模、评估闭环和性能优化方面仍存在提升空间。
\end{abstractcn}
